\chapter{Półmartyngały i wzór Itô dla procesów ze skokami}
\label{ch:semimartingale-jumps}

\section{Dlaczego połączenie części ciągłej i skokowej wymaga nowego języka?}

Wcześniej skonstruowaliśmy całkę Itô względem procesu Wienera,
a w Rozdziale~\ref{ch:poisson-random-measure} --- względem
skompensowanej miary Poissona. Oba przypadki łączy rozkład
zmienności procesu na część lokalnie martyngałową i część o
skończonej wariacji. Ta obserwacja motywuje klasę
\emph{półmartyngałów}. Wprowadzimy ją na poziomie potrzebnym do
jawnych SDE ze skokami, a głębsze twierdzenia o wszystkich
półmartyngałach sformułujemy z wyraźnym zaznaczeniem zakresu dowodu.

\begin{definicja}[Półmartyngał i półmartyngał szczególny]
\label{def:semimartingale}
Niech $(\calF_t)$ będzie filtracją spełniającą warunki zwykłe,
a $X$ adaptowanym procesem càdlàg na $[0,T]$.
Proces $X$ nazywamy \emph{półmartyngałem}, jeżeli istnieją
càdlàg lokalny martyngał $M$ oraz adaptowany càdlàg proces
$A$ o skończonej wariacji na każdym przedziale $[0,t]$,
takie że $X=X_0+M+A$ i $M_0=A_0=0$.
Jeżeli można wybrać $A$ przewidywalny, to $X$ nazywamy
\emph{półmartyngałem szczególnym}. W takim przypadku
rozkład na przewidywalny $A$ i lokalny martyngał $M$
jest jednoznaczny przy powyższej normalizacji.
\end{definicja}

\begin{proof}[Uzasadnienie jednoznaczności rozkładu szczególnego]
Załóżmy, że $X=X_0+M+A=X_0+M'+A'$, gdzie $A,A'$ są
przewidywalne i mają lokalnie skończoną wariację.
Proces $L=M-M'=A'-A$ jest przewidywalnym lokalnym
martyngałem o skończonej wariacji. Ponieważ jego skoki są
przewidywalne, dla każdego czasu przewidywalnego $\tau$
warunek martyngałowy daje
$\E[\Delta L_\tau\mid\calF_{\tau-}]=0$ po lokalizacji.
Ale $\Delta L_\tau$ jest $\calF_{\tau-}$-mierzalny,
a więc $\Delta L_\tau=0$. Twierdzenie o anonsowaniu
skoków procesu przewidywalnego pozwala wywnioskować,
że $L$ nie ma skoków. Ciągły lokalny martyngał o
skończonej wariacji jest stały (jego wariacja kwadratowa
wynosi zero), a $L_0=0$. Otrzymujemy $L=0$.
\end{proof}

\begin{uwaga}[Granica powyższego argumentu]
Użyte w ostatnim dowodzie twierdzenie o anonsowaniu skoków
przewidywalnych jest standardowym wynikiem ogólnej teorii
procesów; jego pełna konstrukcja nie została tutaj
przeprowadzona. W dalszych konkretnych przykładach
jednoznaczność kompensatorów wynika bezpośrednio z
obliczeń warunkowych. Nie zastępujemy w ten sposób
pełnego dowodu ogólnego twierdzenia Dooba--Meyera.
\end{uwaga}

\begin{przyklad}[Trzy różne półmartyngały]
\begin{enumerate}
\item $X_t=W_t+\int_0^tb_s\,\dd s$: część
lokalnie martyngałowa jest ciągła, a skończona wariacja
pochodzi z dryfu.
\item $X_t=N_t-\lambda t$ jest martyngałem, mimo że
między skokami ma ujemne nachylenie; jego wariacja
całkowita jest skończona na każdym zwartym przedziale.
\item Proces gamma jest niemalejącym procesem o
skończonej wariacji, mimo iż na każdym przedziale
niezerowej długości wykonuje nieskończenie wiele skoków.
\end{enumerate}
Drugi przykład pokazuje, że proces o skończonej wariacji
\emph{nie musi} być przewidywalny ani wolny od części
martyngałowej. Sam wygląd trajektorii nie wyznacza
kanonicznej dekompozycji.
\end{przyklad}

\section{Całka względem półmartyngału: konstrukcja w dostępnej klasie}

Zaczynamy od klasy, którą potrafimy zbudować przy pomocy
już udowodnionych izometrii. Niech $X=X_0+M+A$, gdzie
$M$ jest sumą całek względem procesu Wienera i
skompensowanej miary Poissona, lokalnie kwadratowo
całkowalnych w sensie poprzednich rozdziałów, a $A$
adaptowanym procesem càdlàg o lokalnie skończonej wariacji.
Dla ograniczonego procesu przewidywalnego $H$, dla którego
oba składniki $H$ razy funkcje podcałkowe $M$ należą lokalnie
do odpowiednich przestrzeni $L^2$, definiujemy
\[
\int_0^t H_s\,\dd X_s
:=\int_0^tH_s\,\dd M_s+\int_{(0,t]}H_s\,\dd A_s.
\tag{12.1}\label{eq:semi-integral}
\]
Pierwszą całkę definiujemy przez lokalizację w $L^2$,
a drugą trajektoria po trajektorii jako całkę Lebesgue'a--Stieltjesa.
W miejscu skoku współczynnik ma wartość $H_s$, dostępną
z informacji sprzed skoku w sensie przewidywalności. Jeśli
$H_s=Y_{s-}$ dla adaptowanego procesu càdlàg $Y$, jest to
konkretnie lewostronna granica $Y$; nie każda przewidywalna
funkcja podcałkowa musi mieć postać lewostronnej granicy.

\begin{lemat}[Zatrzymywanie i skok całki]
\label{lem:semi-jumps}
W sytuacji z \eqref{eq:semi-integral}, jeśli $H$ jest
ograniczony i przewidywalny, a $\tau\le T$ czasem zatrzymania,
to
\[
\left(\int_0^\cdot H_s\,\dd X_s\right)^\tau_t
=\int_0^tH_s\ind{\{s\le\tau\}}\,\dd X_s,
\qquad
\Delta\left(\int_0^\cdot H_s\,\dd X_s\right)_t
=H_t\Delta X_t.
\]
\end{lemat}
\begin{proof}
Dla całki względem $A$ oba wzory wynikają bezpośrednio
z definicji całki Stieltjesa z atomami i z własności
miary $\dd A$. Dla elementarnego przewidywalnego $H$
są to tożsamości sum przyrostów lokalnego martyngału.
W ogólnym przypadku wykorzystujemy czasy lokalizacji z konstrukcji
całek Wienera i Poissona, aby odpowiednie całki miały skończone
drugie momenty. Aproksymacja
procesów elementarnych i nierówność maksymalna dla
całek martyngałowych dają zbieżność ucp, zatem
przechodzimy do granicy w pierwszej równości.
Jeśli $Y_n\to Y$ ucp i wszystkie procesy są càdlàg,
to wzdłuż podciągu zachodzi zbieżność jednostajna
prawie na pewno, a stąd $\Delta Y_n(t)\to\Delta Y(t)$
jednostajnie w $t$ na ustalonym zwartym przedziale.
To przenosi równość o skokach, a lokalizacja kończy dowód.
\end{proof}

\begin{uwaga}[Co zostaje poza konstrukcją?]
Dla lokalnych martyngałów, które nie są lokalnie
kwadratowo całkowalne, pełne zdefiniowanie
$H\cdot X$ wymaga ogólnej teorii całkowania względem
półmartyngałów. Jej ścisła wersja obejmuje ciągłość
operatora całkowego w topologii półmartyngałów.
Nie utożsamiamy tego rozszerzenia z samą izometrią $L^2$;
w dalszych dowodach pracujemy w jawnej klasie
procesów Wienera i skompensowanych miar Poissona.
\end{uwaga}

\section{Wariacja kwadratowa: część ciągła i rzeczywiste skoki}

\begin{twierdzenie}[Wariacja kwadratowa w klasie dyfuzji ze skokami]
\label{th:jump-qv}
Niech $X$ będzie rzeczywistym procesem postaci
\[
X_t=x+\int_0^t b_s\,\dd s+\int_0^t\sigma_s\,\dd W_s
+\int_{(0,t]\times E_0}\gamma_s(z)\,\widetilde N(\dd s,\dd z)
+\int_{(0,t]\times E_1}\Gamma_s(z)\,N(\dd s,\dd z),
\]
gdzie $E_0,E_1\subset E$ są rozłączne,
$\nu(E_1)<\infty$, a współczynniki są przewidywalne i spełniają
$\int_0^T(|b_s|+\sigma_s^2+\int_{E_0}\gamma_s(z)^2\dd\nu(z))\dd s
<\infty$ prawie na pewno.
Wówczas dla podziałów deterministycznych o średnicy dążącej
do zera sumy kwadratów przyrostów zbiegają według
prawdopodobieństwa do
\[
[X]_t=\int_0^t\sigma_s^2\,\dd s
+\int_{(0,t]\times E_0}\gamma_s(z)^2\,N(\dd s,\dd z)
+\int_{(0,t]\times E_1}\Gamma_s(z)^2\,N(\dd s,\dd z).
\tag{12.2}\label{eq:jump-qv}
\]
Zbieżność dotyczy każdej ustalonej chwili $t$; przy
silniejszych oszacowaniach po lokalizacji otrzymujemy ucp.
\end{twierdzenie}
\begin{proof}
Najpierw obcinamy zakres małych skoków tak, aby jego
intensywność była skończona. Dla procesu mającego
skończenie wiele skoków na $[0,T]$ oraz ciągłą część
Itô możemy rozbić sumę kwadratów przyrostów na
wyrazy: kwadrat ciągłej części, kwadrat skokowej
części i wyrazy mieszane. Dla ciągłej części
wariacja kwadratowa wynosi $\int\sigma^2\dd s$;
ciągła część skończonej wariacji ma wariację kwadratową
zero. Każdy ustalony skok znajdzie się dla dostatecznie
drobnego podziału w odrębnym przedziale, a ciągłość
pozostałej części oznacza, że wyrazy mieszane dążą do
zera. Otrzymujemy \eqref{eq:jump-qv} dla obciętego procesu.

Dla części o nieskończonej aktywności stosujemy
aproksymację $X^{(n)}$ przez ucięcie miary skoków.
Najpierw lokalizujemy procesy tak, aby całkowita wariacja
składnika dryfu oraz całki kwadratów współczynników
ciągłych i małoskokowych były ograniczone stałą. Ponieważ
dużych skoków jest na $[0,T]$ skończenie wiele i ich rozmiary
są skończone prawie na pewno, możemy także obciąć ich
rozmiary i liczbę; prawdopodobieństwo zmiany pierwotnej
trajektorii przez to obcięcie dąży do zera. W tym
ograniczonym modelu izometria Poissona i nierówność maksymalna
dają $X^{(n)}\to X$ ucp oraz
$\E\int\!|(\gamma^{(n)})^2-\gamma^2|\,\dd N\to0$.
Sama zbieżność ucp nie wystarczyłaby do zbieżności wariacji
kwadratowej. Aby uzasadnić przejście dla sum kwadratów
przyrostów, używamy nierówności Cauchy'ego--Schwarza
\[
\left|\sum_j(\Delta_jX)^2-\sum_j(\Delta_jX^{(n)})^2\right|
\le\left(\sum_j(\Delta_j(X-X^{(n)}))^2\right)^{1/2}
\left(\sum_j(\Delta_j(X+X^{(n)}))^2\right)^{1/2}.
\]
Dla sum martyngałowych po lokalizacji pierwszy czynnik
ma drugi moment kontrolowany przez izometrię, jednostajnie
po podziałach, a drugi ma ograniczony drugi moment
(np. przez zatrzymanie procesów i oszacowania dla części
skończonej wariacji). Zatem najpierw wybieramy $n$ duże,
a potem średnicę podziału małą. Otrzymujemy zbieżność
według prawdopodobieństwa do prawej strony
\eqref{eq:jump-qv}. Dla nieograniczonych dużych skoków
stosujemy lokalizację według liczby i rozmiaru skoków.
\end{proof}

\begin{przyklad}[Skoki nie tworzą członu $\dd t$ w chwili skoku]
Dla $M_t=N_t-\lambda t$ zachodzi $[M]_t=N_t$,
a nie $\lambda t$. Proces $\lambda t$ jest jego nawiasem
przewidywalnym. Natomiast dla niezależnego Wienera
$W$ mamy $[W+M]_t=t+N_t$.
\end{przyklad}

\section{Wzór Itô ze skokami: najpierw przypadek skończony}

\begin{lemat}[Wzór dla skończonej liczby skoków]
\label{lem:finite-jump-ito}
Niech $X=C+J$, gdzie $C$ jest ciągłym procesem Itô
na $[0,T]$, a $J_t=\sum_{k:S_k\le t}U_k$ adaptowanym procesem
o skończenie wielu skokach na $[0,T]$ (bez dodatkowej
ciągłej części skończonej wariacji).
Dla $f\in C^{1,2}([0,T]\times\R)$ zachodzi
\[
\begin{split}
f(t,X_t)-f(0,X_0)={}&\int_0^t f_t(s,X_{s-})\,\dd s
+\int_0^t f_x(s,X_{s-})\,\dd X_s\\
&+\frac12\int_0^t f_{xx}(s,X_{s-})\,\dd[C]_s\\
&+\sum_{0<s\le t}\bigl(f(s,X_s)-f(s,X_{s-})
-f_x(s,X_{s-})\Delta X_s\bigr).
\end{split}
\tag{12.3}\label{eq:ito-finite-jumps}
\]
\end{lemat}
\begin{proof}
Uporządkujmy chwile skoków $J$ w $[0,t]$:
$0<\tau_1<\cdots<\tau_m\le t$. Pomiędzy skokami
proces $X$ jest ciągłym procesem Itô z dodaną
ciągłą częścią procesu o skończonej wariacji.
Zastosowanie ciągłego wzoru Itô na kolejnych przedziałach
$[0,\tau_1)$, $(\tau_1,\tau_2)$, itd., po lokalizacji i
przejściu do granic jednostronnych, daje składniki
całkowe. Przy każdym skoku występuje dokładna różnica
$f(\tau_j,X_{\tau_j})-f(\tau_j,X_{\tau_j-})$.
W całce $\int f_x\,\dd X$ skok uwzględniony jest jako
$f_x(\tau_j,X_{\tau_j-})\Delta X_{\tau_j}$.
Ich różnica to właśnie składnik sumy w tezie.
Sumowanie skończonej liczby przedziałów i skoków
kończy dowód. Nie ma tu przejścia do granicy w liczbie skoków.
\end{proof}

\section{Wzór Itô dla kompensowanych małych skoków}

\begin{twierdzenie}[Wzór Itô dla skokowej dyfuzji]
\label{th:ito-jump-diffusion}
Niech $T>0$, $N$ będzie losową miarą Poissona na
$(0,T]\times E$ o intensywności $\dd s\,\dd\nu(z)$,
$W$ zaś procesem Wienera względem wspólnej filtracji,
niezależnym od $N$. Niech $E=E_0\sqcup E_1$ i
$\nu(E_1)<\infty$. Niech $b,\sigma$ będą rzeczywistymi
procesami przewidywalnymi, a $\gamma,\Gamma$ funkcjami
przewidywalnymi na $\Omega\times(0,T]\times E_0$ i $E_1$
odpowiednio, spełniającymi
\[
\E\int_0^T\left(|b_s|+\sigma_s^2+
\int_{E_0}\gamma_s(z)^2\,\dd\nu(z)\right)\dd s<\infty.
\]
Nie nakładamy warunku momentowego na duże skoki $\Gamma$;
zakładamy jedynie jej wartości rzeczywiste $N$-prawie wszędzie.
Niech
\[
X_t=x+\int_0^tb_s\,\dd s+\int_0^t\sigma_s\,\dd W_s
+\int_{(0,t]\times E_0}\gamma_s(z)\,\widetilde N(\dd s,\dd z)
+\int_{(0,t]\times E_1}\Gamma_s(z)\,N(\dd s,\dd z).
\]
Dla $f\in C^{1,2}([0,T]\times\R)$, którego
$f,f_t,f_x,f_{xx}$ są ograniczone, zachodzi prawie na pewno
jednocześnie dla wszystkich $t\in[0,T]$:
\[
\begin{split}
f(t,X_t)={}&f(0,x)
+\int_0^t\left(f_t+b_sf_x+\tfrac12\sigma_s^2f_{xx}\right)(s,X_{s-})\,\dd s\\
&+\int_0^t\sigma_s f_x(s,X_{s-})\,\dd W_s\\
&+\int_{(0,t]\times E_0}
\bigl[f(s,X_{s-}+\gamma_s(z))-f(s,X_{s-})\bigr]
\,\widetilde N(\dd s,\dd z)\\
&+\int_0^t\int_{E_0}
\bigl[f(s,X_{s-}+\gamma_s(z))-f(s,X_{s-})
-f_x(s,X_{s-})\gamma_s(z)\bigr] \,\dd\nu(z)\,\dd s\\
&+\int_{(0,t]\times E_1}
\bigl[f(s,X_{s-}+\Gamma_s(z))-f(s,X_{s-})\bigr]
\,N(\dd s,\dd z).
\end{split}
\tag{12.4}\label{eq:ito-jump-general}
\]
We wszystkich wyrażeniach $f_t,f_x,f_{xx}$ są oceniane
w punkcie $(s,X_{s-})$, o ile nie zapisano inaczej.
\end{twierdzenie}
\begin{proof}
\emph{Krok 1: obcięcie małych skoków.} Wybieramy rosnące
zbiory $E_0^{(n)}\uparrow E_0$ z
$\nu(E_0^{(n)})<\infty$ i
$\E\int_0^T\int_{E_0\setminus E_0^{(n)}}\gamma_s(z)^2
\,\dd\nu(z)\,\dd s\to0$ (istnieją przez $\sigma$-skończoność
$\nu$ oraz zbieżność zdominowaną).
Zastępujemy w definicji $X$ małą skompensowaną całkę
przez całkę po $E_0^{(n)}$, otrzymując $X^{(n)}$.
Z izometrii \eqref{eq:prm-isometry} i nierówności Dooba
\[
\E\sup_{t\le T}|X_t-X_t^{(n)}|^2
\le4\E\int_0^T\int_{E_0\setminus E_0^{(n)}}\gamma_s(z)^2
\,\dd\nu(z)\,\dd s\longrightarrow0.
\]

\emph{Krok 2: skończenie wiele skoków.} Na $E_0^{(n)}$
całka skompensowana ma postać sumy skoków minus
ciągły proces skończonej wariacji
$\int_0^t\int_{E_0^{(n)}}\gamma_s(z)\,\dd\nu(z)\,\dd s$.
Ten drugi proces ma skończoną wariację prawie na pewno,
ponieważ Cauchy--Schwarz i skończoność $\nu(E_0^{(n)})$
dają całkowalność bezwzględną. Wobec tego
Lemat~\ref{lem:finite-jump-ito} stosuje się do $X^{(n)}$.
Rozpisujemy wszystkie skoki i grupujemy ich wkłady:
całka $f_x\gamma\,\widetilde N$ oraz suma poprawek
$R=f(X_-+\gamma)-f(X_-)-f_x(X_-)\gamma$ dają
po przekształceniu
$[f(X_-+\gamma)-f(X_-)]\,\widetilde N+R\,\dd s\,\dd\nu$.
Duże skoki pozostają nieskompensowane. Otrzymujemy
\eqref{eq:ito-jump-general} z $E_0$ zastąpionym przez
$E_0^{(n)}$ i $X$ przez $X^{(n)}$.

\emph{Krok 3: przejście do granicy w członach ciągłych.}
Przechodzimy do podciągu, po którym $X^{(n)}\to X$
jednostajnie prawie na pewno. Na każdym przedziale
$[0,T]$ trajektoria $X$ jest ograniczona; jednostajna
ciągłość pochodnych $f$ na zwartych prostokątach
i zbieżność zdominowana prowadzą do zbieżności
zwykłych całek po czasie. Dla całki Wienera stosujemy
izometrię i ograniczenie $|f_x|\le C$:
\[
\E\int_0^T\sigma_s^2
|f_x(s,X^{(n)}_{s-})-f_x(s,X_{s-})|^2\,\dd s\to0.
\]
Wynika to z dominacji przez $4C^2\sigma_s^2$ oraz
zbieżności punktowej po podciągu, prawie wszędzie
w $\Prob\otimes\dd s$.

\emph{Krok 4: małe skoki.} Z twierdzenia o wartości
średniej mamy
$|f(s,x+u)-f(s,x)|\le C|u|$ i
$|R(s,x,u)|\le\frac12C u^2$, gdzie $C$ ogranicza
$|f_x|,|f_{xx}|$. Izometria Poissona pozwala więc
przejść do granicy w skompensowanych całkach:
dla wspólnej części $E_0^{(n)}$ różnice argumentów
zbiegają punktowo i są dominowane przez $2C|\gamma|$,
a na usuniętej części $E_0\setminus E_0^{(n)}$
norma kwadratowa dąży do zera. Podobnie poprawki
$R(s,X^{(n)}_{s-},\gamma_s(z))$ zbiegają w
$L^1(\Prob\otimes\dd s\otimes\nu)$ dzięki dominacji
przez $C\gamma^2$. Po lokalizacji pozostaje też
ciągłość składnika dużych skoków: na $[0,T]$ jest ich
skończenie wiele, a argumenty $X^{(n)}_{s-}$ zbiegają
jednostajnie wzdłuż podciągu. Wszystkie człony prawej
strony zbiegają zatem do odpowiednich członów w tezie.
Zbieżność ucp i dobór wspólnego podciągu dla skończonej
rodziny całek zapewniają równość procesów z dokładnością
do nierozróżnialności.
\end{proof}

\begin{uwaga}[Dlaczego poprawka ma rząd kwadratowy?]
Zapis $f(x+u)-f(x)-f'(x)u$ jest błędem liniowego
przybliżenia funkcji. Przy $f\in C^2$ jego wielkość jest
rzędu $u^2$; dlatego właśnie warunek
$\int_{|z|\le1}|z|^2\,\dd\nu(z)<\infty$
wystarcza do całkowania poprawek po nieskończenie wielu
małych skokach. To jest odpowiednik roli wariacji
kwadratowej ruchu Browna w ciągłym wzorze Itô.
\end{uwaga}

\begin{zadanie}[Dlaczego zwykły wzór łańcuchowy jest niepoprawny?]
Niech $N$ będzie procesem Poissona o intensywności
$\lambda>0$ i $f(x)=x^2$. Wyznaczyć $N_t^2$ przez
sumowanie po skokach oraz wskazać brakujący człon
w naiwnym wzorze $\dd f(N_t)=f'(N_{t-})\,\dd N_t$.

\emph{Rozwiązanie.} W każdej chwili skoku $\Delta N_s=1$,
więc $(N_{s-}+1)^2-N_{s-}^2=2N_{s-}+1$.
Sumując po chwilach skoków, otrzymujemy
\[
N_t^2=\int_{(0,t]}2N_{s-}\,\dd N_s+N_t.
\]
Drugi składnik jest dokładnie
$\sum_{s\le t}(\Delta N_s)^2$.
Gdybyśmy bez poprawki zastosowali klasyczną
regułę łańcuchową, zgubilibyśmy jeden przy każdym skoku.
Dla gładkiego $f$ odpowiednia poprawka jest równa
$f(N_{s-}+1)-f(N_{s-})-f'(N_{s-})$.
\end{zadanie}

\section{Iloczyny, wykładnik stochastyczny i zastosowania}

\begin{wniosek}[Reguła iloczynu z kowariacją skokową]
\label{cor:jump-product}
Dla dwóch procesów $X,Y$ klasy z
Twierdzenia~\ref{th:ito-jump-diffusion}, napędzanych
wspólnym procesem Wienera i wspólną miarą Poissona,
zachodzi
\[
X_tY_t=X_0Y_0+\int_0^tX_{s-}\,\dd Y_s
+\int_0^tY_{s-}\,\dd X_s+[X,Y]_t,
\]
gdzie $[X,Y]$ zawiera kowariację części ciągłych
oraz sumę $\sum_{s\le t}\Delta X_s\Delta Y_s$.
\end{wniosek}
\begin{proof}
Stosujemy wzór Itô do dwuwymiarowego procesu
$(X,Y)$ i funkcji $f(x,y)=xy$, po uprzednim zatrzymaniu
procesów na przedziałach, na których $|X|+|Y|$ jest
ograniczone. Poza skokami jedyny składnik drugiego
rzędu to $\dd[X,Y]^c$, a na skoku zachodzi dokładna
tożsamość
$(x+u)(y+v)-xy-yu-xv=uv$.
Sumowanie tych składników jest definicją kowariacji.
Formalną wielowymiarową wersję uzyskujemy również
przez polaryzację jednowymiarowego wzoru dla kwadratu.
\end{proof}

\begin{twierdzenie}[Wykładnik stochastyczny w przypadku skończonej aktywności]
\label{th:dolean-finite}
Niech $Z$ będzie rzeczywistym procesem postaci
$Z_t=\int_0^t a_s\,\dd s+\int_0^t\sigma_s\,\dd W_s
+\sum_{k:S_k\le t}U_k$, gdzie na $[0,T]$ suma zawiera
skończenie wiele skoków, $a,\sigma$ są przewidywalne,
$\int_0^T(|a_s|+\sigma_s^2)\dd s<\infty$ prawie na pewno, a
$U_k>-1$. Wówczas równanie
$\dd Y_t=Y_{t-}\,\dd Z_t$, $Y_0=y_0$, ma dokładnie jedno
adaptowane rozwiązanie càdlàg oraz
\[
Y_t=y_0\exp\left(Z_t-Z_0-\frac12\int_0^t\sigma_s^2\,\dd s
-\sum_{S_k\le t}U_k\right)
\prod_{S_k\le t}(1+U_k).
\tag{12.5}\label{eq:dolean-finite}
\]
\end{twierdzenie}
\begin{proof}
Między kolejnymi skokami równanie jest ciągłym liniowym
SDE. Wzór Itô zastosowany do wykładnika pokazuje,
że jego rozwiązanie ma postać
$Y_s=Y_r\exp(Z_s-Z_r-\frac12\int_r^s\sigma_u^2\dd u)$.
W chwili $S_k$ równanie całkowe daje
$\Delta Y_{S_k}=Y_{S_k-}U_k$, a więc
$Y_{S_k}=Y_{S_k-}(1+U_k)$.
Łącząc te równości dla skończonej liczby skoków,
otrzymujemy \eqref{eq:dolean-finite}. Jednoznaczność
wynika z jednoznaczności liniowego SDE między skokami
i jednoznacznego przeliczenia w każdej chwili skoku.
Dodatniość $1+U_k$ zapewnia zachowanie znaku $Y$;
gdy $U_k=-1$, rozwiązanie po skoku staje się zerem,
co wymaga osobnego uwzględnienia we wzorze logarytmicznym.
\end{proof}

\begin{zadanie}[Cena z losowymi skokami i poprawka dryfu]
Niech $\mu\in\R$, $\sigma\ge0$, $\lambda>0$, $\kappa>-1$
oraz $N$ będzie procesem Poissona o intensywności $\lambda$,
niezależnym od $W$. Rozwiązać
\[
\dd S_t=\mu S_{t-}\,\dd t+\sigma S_{t-}\,\dd W_t
+\kappa S_{t-}\,\dd(N_t-\lambda t),\qquad S_0=s_0>0.
\]
Obliczyć $\E S_t$ i wariancję, jeżeli jest skończona.

\emph{Rozwiązanie.} Między skokami efektywny dryf to
$\mu-\lambda\kappa$. Każdy skok mnoży wartość procesu
przez $1+\kappa$. Otrzymujemy
\[
S_t=s_0\exp\left((\mu-\lambda\kappa-\tfrac12\sigma^2)t
+\sigma W_t\right)(1+\kappa)^{N_t}.
\]
Niezależność daje
$\E e^{\sigma W_t-\sigma^2t/2}=1$ oraz
$\E(1+\kappa)^{N_t}=e^{\lambda t\kappa}$,
zatem $\E S_t=s_0e^{\mu t}$.
Dla drugiego momentu używamy
$\E e^{2\sigma W_t-\sigma^2t}=e^{\sigma^2t}$ oraz
$\E(1+\kappa)^{2N_t}=e^{\lambda t(2\kappa+\kappa^2)}$.
Stąd
\[
\E S_t^2=s_0^2e^{(2\mu+\sigma^2+\lambda\kappa^2)t},
\quad
\Var(S_t)=s_0^2e^{2\mu t}
\bigl(e^{(\sigma^2+\lambda\kappa^2)t}-1\bigr).
\]
To zadanie pokazuje, dlaczego kompensacja zmienia dryf,
ale nie rozmiar pojedynczego skoku.
\end{zadanie}

\begin{zadanie}[Proces Ornsteina--Uhlenbecka wymuszany skokami]
Niech $N$ będzie procesem Poissona o intensywności
$\lambda>0$, $a>0$, a $J_t=c(N_t-\lambda t)$ dla
$c\in\R$. Rozwiązać równanie
$\dd X_t=-aX_{t-}\dd t+\dd J_t$ z deterministycznym
$X_0=x$ oraz obliczyć $\E X_t$ i $\Var(X_t)$.

\emph{Rozwiązanie.} Całkowanie przez części dla
$e^{at}X_t$ (czynnik deterministyczny nie ma skoków)
daje
$X_t=xe^{-at}+c\int_0^te^{-a(t-s)}\,\dd(N_s-\lambda s)$.
Całka ma średnią zero, więc $\E X_t=xe^{-at}$.
Izometria z Rozdziału~\ref{ch:poisson-random-measure} daje
\[
\Var(X_t)=\lambda c^2\int_0^t e^{-2a(t-s)}\,\dd s
=\frac{\lambda c^2}{2a}(1-e^{-2at}).
\]
W odróżnieniu od gaussowskiego procesu
Ornsteina--Uhlenbecka rozkład $X_t$ ma tu, dla
$c\ne0$, składnik atomowy odpowiadający brakowi skoków
na $[0,t]$. Nie jest zatem ogólnie rozkładem normalnym.
\end{zadanie}

\section{Generator dyfuzji ze skokami i równanie wsteczne}

\begin{twierdzenie}[Generator w przypadku jednorodnych współczynników]
\label{th:jump-generator}
Niech $X^x$ będzie czasowo jednorodnym procesem Markowa
spełniającym równanie skokowe
z $b=b(x)$, $\sigma=\sigma(x)$, $\gamma=\gamma(x,z)$,
$\Gamma=\Gamma(x,z)$, z istnieniem rozwiązania i
odpowiednią ciągłością w warunku początkowym.
Dla $f\in C_b^2(\R)$ niech
\[
\begin{split}
Lf(x)={}&b(x)f'(x)+\frac12\sigma(x)^2f''(x)\\
&+\int_{E_0}[f(x+\gamma(x,z))-f(x)-f'(x)\gamma(x,z)]\,\dd\nu(z)\\
&+\int_{E_1}[f(x+\Gamma(x,z))-f(x)]\,\dd\nu(z).
\end{split}
\tag{12.6}\label{eq:jump-generator}
\]
Jeżeli powyższe całki są skończone i $Lf$ jest
ograniczona oraz ciągła, to
$f(X_t^x)-f(x)-\int_0^tLf(X_{s-}^x)\,\dd s$
jest lokalnym martyngałem. Ponadto, gdy można
uzasadnić zbieżność zdominowaną dla $Lf(X_s^x)$ przy
$s\downarrow0$, zachodzi
$\lim_{t\downarrow0}t^{-1}(\E f(X_t^x)-f(x))=Lf(x)$.
\end{twierdzenie}
\begin{proof}
Wstawiamy do wzoru \eqref{eq:ito-jump-general}
$f(t,x)=f(x)$. Dla dużych skoków, których liczba jest
skończona na odcinkach zwartych, formuła kompensacji
\eqref{eq:compensation} pokazuje, że
\[
\int_{(0,t]\times E_1}\Delta_zf(X_{s-})\,N(\dd s,\dd z)
-\int_0^t\int_{E_1}\Delta_zf(X_{s-})\,\dd\nu(z)\,\dd s
\]
jest lokalnym martyngałem; $\Delta_zf(x)=f(x+\Gamma(x,z))-f(x)$.
Składniki Wienera i małych skoków są lokalnymi martyngałami
z poprzednich rozdziałów. Pozostałe człony skończonej
wariacji sumują się dokładnie do $\int Lf(X_{s-})\dd s$.
Po wzięciu wartości oczekiwanych (przy zapewnionej
całkowalności) otrzymujemy
$\E f(X_t^x)-f(x)=\E\int_0^tLf(X_s^x)\dd s$.
Ciągłość stochastyczna w zerze i ciągłość $Lf$ prowadzą
do $\E Lf(X_s^x)\to Lf(x)$ przy $s\downarrow0$,
a ograniczoność $Lf$ pozwala zastosować zbieżność
zdominowaną. Dzielenie przez $t$ kończy dowód.
\end{proof}

\begin{przyklad}[Generator procesu złożonego Poissona]
Dla $X_t^x=x+\sum_{j=1}^{N_t}Y_j$, gdzie
$N$ ma intensywność $\lambda$ i znaczniki mają prawo $\rho$,
$Lf(x)=\lambda\int_{\R}(f(x+y)-f(x))\,\dd\rho(y)$.
Jeśli $Y_j=a$ deterministycznie, mamy
$Lf(x)=\lambda(f(x+a)-f(x))$ i równanie wsteczne
$\partial_tu(t,x)=\lambda[u(t,x+a)-u(t,x)]$.
Dla $u(0,x)=f(x)$ jego rozwiązanie dane przez
półgrupę ma postać
$u(t,x)=\sum_{n\ge0}e^{-\lambda t}(\lambda t)^n f(x+na)/n!$.
Sprawdzenie różniczkowania szeregu jest dozwolone dla
ograniczonej $f$, bo szereg pochodnych po $t$ zbiega
jednostajnie na zwartych przedziałach czasu.
\end{przyklad}

\begin{center}
\begin{tikzpicture}[>=Latex,scale=.95]
\draw[->] (0,0)--(9,0) node[right] {$t$};
\draw[->] (0,-.8)--(0,3.0) node[above] {$X_t$};
\draw[blue!75!black,thick] plot[smooth] coordinates
{(0,.5) (.8,.7) (1.5,.45) (2.3,.8) (3,1.0) (3.8,.65)};
\draw[blue!75!black,thick] (3.8,.65)--(3.8,1.85);
\draw[blue!75!black,thick] plot[smooth] coordinates
{(3.8,1.85) (4.5,2.02) (5.2,1.75) (5.8,2.0)};
\draw[blue!75!black,thick] (5.8,2.0)--(5.8,1.1);
\draw[blue!75!black,thick] plot[smooth] coordinates
{(5.8,1.1) (6.7,1.5) (7.5,1.2) (8.5,1.65)};
\draw[orange!80!black,dashed,thick] (3.8,.65)--(3.8,1.85);
\draw[orange!80!black,dashed,thick] (5.8,1.1)--(5.8,2.0);
\node[anchor=west,scale=.87,orange!80!black] at (6.1,2.55)
{skoki $\Delta X$};
\node[anchor=west,scale=.87,blue!70!black] at (.7,2.52)
{część ciągła};
\end{tikzpicture}
\par\smallskip\small Trajektoria skokowej dyfuzji:
między skokami zachodzi ciągła ewolucja, a w chwilach
losowych pojawiają się nieciągłości. Wzór Itô uwzględnia
obie części i dokładną poprawkę $f(X_s)-f(X_{s-})$.
\end{center}

\section{Co zostało dowiedzione, a co należy do ogólnej teorii?}

W tym rozdziale udowodniliśmy wzór Itô dla
\emph{dyfuzji z małymi skokami całkowanymi w $L^2$
i skończoną intensywnością dużych skoków}.
To wystarcza do całych rodzin procesów Lévy'ego oraz
wielu równań skokowych. Nie dowiedliśmy natomiast
ogólnego twierdzenia charakteryzującego wszystkie
półmartyngały jako dobre całkowniki, ani ogólnego
rozkładu Dooba--Meyera dla podmartyngałów klasy (D).
Podajemy ich właściwe miejsce w teorii, ale nie sugerujemy,
że wynikają z izometrii Poissona bez dalszej pracy.
W kolejnej części tomu zastosujemy przedstawiony rachunek
do równań stochastycznych ze skokami, a następnie wrócimy
do ścisłych całek po trajektoriach i mechaniki kwantowej.

% ===========================================================================
% ROZDZIAŁ 13 -- Stochastyczne równania różniczkowe ze skokami
% ===========================================================================
